\chapter{Einleitung}
\todo{Auf Smarthome integration noch im Fazit eingehen}
\todo{Abstract deutsch überarbeiten}
\todo{Abstract englisch übersetzen}
\todo{Abbildung und Abb. Checken damit einheitlich}
\todo{alle cites kontrollier}
\todo{alle werte mit Einheiten und ~ kontrollieren}
\todo{Absätze prüfen ob zusammenführen sinnvoll}
\todo{Tablett korrektur einarbeiten}
Diese Arbeit befasst sich mit der Konzeptionierung, Entwicklung und dem Bau einer so benannten WordClock mit zusätzlichen informativen Funktionen. 
Dabei werden sowohl technische als auch konstruktive und gestalterische Aspekte betrachtet. 

Die Darstellung der Uhrzeit stellt eine grundlegende Funktion technischer Systeme im Alltag dar. 
Neben analogen und digitalen Anzeigen existieren alternative Konzepte der Zeitdarstellung. 
Eine WordClock stellt die Uhrzeit in sprachlicher Form dar. 
Dabei werden einzelne Wörter gezielt beleuchtet, so dass eine lesbare Zeitangabe entsteht. 

Im Bereich der Elektrotechnik vereint eine WordClock mehrere Fachdisziplinen. 
Hierzu zählen insbesondere Mikrocontrollertechnik, digitale Signalverarbeitung die Integration externer Datenquellen. 
Zusätzlich spielen Aspekte wie Vernetzung, Benutzerinteraktion und Systemintegration eine Rolle. 

Die Entwicklung technischer Systeme erfordert ein strukturiertes Vorgehen. 
Neben der funktionalen Auslegung sind auch Anforderungen an Zuverlässigkeit, Energieeffizienz, Erweiterbarkeit und Gestaltung zu berücksichtigen. 
Eine systematische Konzeption bildet daher die Grundlage für eine zielgerichtete Umsetzung.